\chapter{The Audience}

\lettrine[ante=`]{W}{ell?}' cried Athos with a mild look of reproach, when D'Artagnan had read the letter addressed to him by Monk.

\zz
<Well!> said D'Artagnan, red with pleasure, and a little with shame, at having so hastily accused the king and Monk. <This is a politeness,—which leads to nothing, it is true, but yet it is a politeness.>

<I had great difficulty in believing the young prince ungrateful,> said Athos.

<The fact is, that his present is still too near his past,> replied D'Artagnan; <after all, everything to the present moment proved me right.>

<I acknowledge it, my dear friend, I acknowledge it. Ah! there is your cheerful look returned. You cannot think how delighted I am.>

<Thus you see,> said D'Artagnan, <\regnal{Charles II} receives M.~Monk at nine o'clock; he will receive me at ten; it is a grand audience, of the sort which at the Louvre are called <distributions of court holy water.> Come, let us go and place ourselves under the spout, my dear friend! Come along.>

Athos replied nothing; and both directed their steps, at a quick pace, towards the palace of St James's, which the crowd still surrounded, to catch, through the windows, the shadows of the courtiers, and the reflection of the royal person. Eight o'clock was striking when the two friends took their places in the gallery filled with courtiers and politicians. Every one looked at these simply-dressed men in foreign costumes, at these two noble heads so full of character and meaning. On their side, Athos and D'Artagnan, having with two glances taken the measure of the whole assembly, resumed their chat.

A great noise was suddenly heard at the extremity of the gallery,—it was General Monk, who entered, followed by more than twenty officers, all eager for a smile, as only the evening before he was master of all England, and a glorious to-morrow was looked to, for the restorer of the Stuart family.

<Gentlemen,> said Monk, turning round, <henceforward I beg you to remember that I am no longer anything. Lately I commanded the principal army of the republic; now that army is the king's, into whose hands I am about to surrender, at his command, my power of yesterday.>

Great surprise was painted on all the countenances, and the circle of adulators and suppliants which surrounded Monk an instant before, was enlarged by degrees, and ended by being lost in the large undulations of the crowd. Monk was going into the ante-chamber as others did. D'Artagnan could not help remarking this to the Comte de la Fère, who frowned on beholding it. Suddenly the door of the royal apartment opened, and the young king appeared, preceded by two officers of his household.

<Good evening, gentlemen,> said he. <Is General Monk here?>

<I am here, sire,> replied the old general.

Charles stepped hastily towards him, and seized his hand with the warmest demonstration of friendship. <General,> said the king, aloud, <I have just signed your patent,—you are Duke of Albemarle; and my intention is that no one shall equal you in power and fortune in this kingdom, where—the noble Montrose excepted—no one has equalled you in loyalty, courage, and talent. Gentlemen, the duke is commander of our armies of land and sea; pay him your respects, if you please, in that character.>

Whilst every one was pressing round the general, who received all this homage without losing his impassibility for an instant, D'Artagnan said to Athos: <When one thinks that this duchy, this commander of the land and sea forces, all these grandeurs, in a word, have been shut up in a box six feet long and three feet wide\longdash>

<My friend,> replied Athos, <much more imposing grandeurs are confined in boxes still smaller,—and remain there forever.>

All at once Monk perceived the two gentlemen, who held themselves aside until the crowd had diminished; he made himself a passage towards them, so that he surprised them in the midst of their philosophical reflections. <Were you speaking of me?> sad he, with a smile.

<My lord,> replied Athos, <we were speaking likewise of God.>

Monk reflected for a moment, and then replied gayly: <Gentlemen, let us speak a little of the king likewise, if you please; for you have, I believe, an audience of his majesty.>

<At nine o'clock,> said Athos.

<At ten o'clock,> said D'Artagnan.

<Let us go into this closet at once,> replied Monk, making a sign to his two companions to precede him; but to that neither would consent.

The king, during this discussion so characteristic of the French, had returned to the centre of the gallery.

<Oh! my Frenchmen!> said he, in that tone of careless gayety which, in spite of so much grief and so many crosses, he had never lost. <My Frenchmen! my consolation!> Athos and D'Artagnan bowed.

<Duke, conduct these gentlemen into my study. I am at your service, messieurs,> added he in French. And he promptly expedited his court, to return to his Frenchmen, as he called them. <Monsieur d'Artagnan,> said he, as he entered his closet, <I am glad to see you again.>

<Sire, my joy is at its height, at having the honour to salute your majesty in your own palace of St James's.>

<Monsieur, you have been willing to render me a great service, and I owe you my gratitude for it. If I did not fear to intrude upon the rights of our command general, I would offer you some post worthy of you near our person.>

<Sire,> replied D'Artagnan, <I have quitted the service of the king of France, making a promise to my prince not to serve any other king.>

<Humph!> said Charles, <I am sorry to hear that; I should like to do much for you; I like you very much.>

<Sire\longdash>

<But, let us see,> said Charles with a smile, <if we cannot make you break your word. Duke, assist me. If you were offered, that is to say, if I offered you the chief command of my musketeers?> D'Artagnan bowed lower than before.

<I should have the regret to refuse what your gracious majesty would offer me,> said he; <a gentleman has but his word, and that word, as I have had the honour to tell your majesty, is engaged to the king of France.>

<We shall say no more about it, then,> said the king, turning towards Athos, and leaving D'Artagnan plunged in the deepest pangs of disappointment.

<Ah! I said so!> muttered the musketeer. <Words! words! Court holy water! Kings have always a marvellous talent for offering us that which they know we will not accept, and in appearing generous without risk. So be it!—triple fool that I was to have hoped for a moment!>

During this time, Charles took the hand of Athos. <Comte,> said he, <you have been to me a second father; the services you have rendered to me are above all price. I have, nevertheless, thought of a recompense. You were created by my father a Knight of the Garter—that is an order which all the kings of Europe cannot bear; by the queen regent, Knight of the Holy Ghost—which is an order not less illustrious; I join to it that of the Golden Fleece sent me by the king of France, to whom the king of Spain, his father-in-law, gave two on the occasion of his marriage; but in return, I have a service to ask of you.>

<Sire,> said Athos, with confusion, <the Golden Fleece for me! when the king of France is the only person in my country who enjoys that distinction?>

<I wish you to be in your country and all others the equal of all those whom sovereigns have honoured with their favour,> said Charles, drawing the chain from his neck; <and I am sure, comte, my father smiles on me from his grave.>

<It is unaccountably strange,> said D'Artagnan to himself, whilst his friend, on his knees, received the eminent order which the king conferred on him—<it is almost incredible that I have always seen showers of prosperity fall upon all who surrounded me, and that not a drop ever reached me! If I were a jealous man, it would be enough to make one tear one's hair, parole d'honneur!>

Athos rose from his knees, and Charles embraced him tenderly. <General!> said he to Monk—then stopping, with a smile, <pardon me, duke, I mean. No wonder if I make a mistake; the word duke is too short for me, I always seek some title to lengthen it. I should wish to see you so near my throne, that I might say to you, as to \regnal{Louis XIV}, my brother! Oh! I have it; and you will almost be my brother, for I make you viceroy of Ireland and Scotland, my dear duke. So, after that fashion, henceforward I shall not make a mistake.>

The duke seized the hand of the king, but without enthusiasm, without joy, as he did everything. His heart, however, had been moved by this last favour. Charles, by skilfully husbanding his generosity, had given the duke time to wish, although he might not have wished for so much as was given him.

<\swear{Mordioux!}> grumbled D'Artagnan, <there is the shower beginning again! Oh! it is enough to turn one's brain!> and he turned away with an air so sorrowful and so comically piteous, that the king, who caught it, could not restrain a smile. Monk was preparing to leave the room, to take leave of Charles.

<What! my trusty and well-beloved!> said the king to the duke, <are you going?>

<With your majesty's permission, for in truth I am weary. The emotions of the day have worn me out; I stand in need of rest.>

<But,> said the king, <you are not going without M.~ d'Artagnan, I hope.>

<Why not, sire?> said the old warrior.

<Well! you know very well why,> said the king.

Monk looked at Charles with astonishment.

<Oh! it may be possible; but if you forget, you, M.~d'Artagnan, do not.>

Astonishment was painted on the face of the musketeer.

<Well, then, duke,> said the king, <do you not lodge with M.~ d'Artagnan?>

<I had the honour of offering M.~d'Artagnan a lodging; yes, sire.>

<That idea is your own, and yours solely?>

<Mine and mine only; yes, sire.>

<Well! but it could not be otherwise—the prisoner always lodges with his conqueror.>

Monk coloured in his turn. <Ah! that is true,> said he; <I am M.~d'Artagnan's prisoner.>

<Without doubt, duke, since you are not yet ransomed; but have no care of that; it was I who took you out of M.~d'Artagnan's hands, and it is I who will pay your ransom.>

The eyes of D'Artagnan regained their gayety and their brilliancy. The Gascon began to understand. Charles advanced towards him.

<The general,> said he, <is not rich, and cannot pay you what he is worth. I am richer, certainly; but now that he is a duke, and if not a king, almost a king, he is worth a sum I could not perhaps pay. Come, M.~ d'Artagnan, be moderate with me; how much do I owe you?>

D'Artagnan, delighted at the turn things were taking, but not for a moment losing his self-possession, replied,—<Sire, your majesty has no occasion to be alarmed. When I had the good fortune to take his grace, M.~Monk was only a general; it is therefore only a general's ransom that is due to me. But if the general will have the kindness to deliver me his sword, I shall consider myself paid; for there is nothing in the world but the general's sword which is worth as much as himself.>

<Odds fish! as my father said,> cried Charles. <That is a gallant proposal, and a gallant man, is he not, duke?>

<Upon my honour, yes, sire,> and he drew his sword. <Monsieur,> said he to D'Artagnan, <here is what you demand. Many have handled a better blade; but however modest mine may be, I have never surrendered it to any one.>

D'Artagnan received with pride the sword which had just made a king.

<Oh! oh!> cried \regnal{Charles II}; <what a sword that has restored me to my throne—to go out of the kingdom—and not, one day, to figure among the crown jewels! No, on my soul! that shall not be! Captain d'Artagnan, I will give you two hundred thousand livres for your sword! If that is too little, say so.>

<It is too little, sire,> replied D'Artagnan, with inimitable seriousness. <In the first place, I do not at all wish to sell it; but your majesty desires me to do so, and that is an order. I obey, then, but the respect I owe to the illustrious warrior who hears me, commands me to estimate a third more the reward of my victory. I ask then three hundred thousand livres for the sword, or I shall give it to your majesty for nothing.> And taking it by the point he presented it to the king. Charles broke into hilarious laughter.

\begin{bwbigpic}
	[\basicscale] 
	{33point}
	[Taking it by the point he presented it to the king]
\end{bwbigpic}

<A gallant man, and a merry companion! Odds fish! is he not, duke? is he not, comte? He pleases me! I like him! Here, Chevalier d'Artagnan, take this.> And going to the table, he took a pen and wrote an order upon his treasurer for three hundred thousand livres.

D'Artagnan took it, and turning gravely towards Monk: <I have still asked too little, I know,> said he, <but believe me, your grace, I would rather have died that allow myself to be governed by avarice.>

The king began to laugh again, like the happiest cockney of his kingdom.

<You will come and see me again before you go, chevalier?> said he; <I shall want to lay in a stock of gayety now my Frenchmen are leaving me.>

<Ah! sire, it will not be with the gayety as with the duke's sword; I will give it to your majesty gratis,> replied D'Artagnan, whose feet scarcely seemed to touch the ground.

<And you, comte,> added Charles, turning towards Athos, <come again, also; I have an important message to confide to you. Your hand, duke.> Monk pressed the hand of the king.

<Adieu! gentlemen,> said Charles, holding out each of his hands to the two Frenchmen, who carried them to their lips.

<Well,> said Athos, when they were out of the palace, <are you satisfied?>

<Hush!> said D'Artagnan, wild with joy, <I have not yet returned from the treasurer's—a shutter may fall upon my head.> 